\documentclass[11pt]{article}

\usepackage[a4paper,margin=0.78in]{geometry}
\usepackage[T1]{fontenc}
\usepackage[utf8]{inputenc}
\usepackage{lmodern}
\usepackage{amsmath,amssymb,amsthm,mathtools}
\usepackage{enumitem,booktabs,array,tabularx,microtype,xcolor}
\usepackage[numbers,sort&compress]{natbib}
\usepackage[colorlinks=true,linkcolor=blue!65!black,
  citecolor=blue!65!black,urlcolor=blue!70!black]{hyperref}
\usepackage[nameinlink,noabbrev]{cleveref}
\setlength{\emergencystretch}{3em}

\hypersetup{
  pdftitle={A Small-Polylogarithmic Determinant-Two Mobius Corridor},
  pdfauthor={Liang Wang},
  pdfsubject={Uniform actual-core cancellation with bounded periodic data},
  pdfkeywords={Mobius correlation, determinant two, exceptional scales, periodic weights}
}

\newtheorem{theorem}{Theorem}[section]
\newtheorem{proposition}[theorem]{Proposition}
\newtheorem{corollary}[theorem]{Corollary}
\newtheorem{lemma}[theorem]{Lemma}
\theoremstyle{definition}
\newtheorem{definition}[theorem]{Definition}
\theoremstyle{remark}
\newtheorem{remark}[theorem]{Remark}

\newcommand{\Lzero}{\mathrm{L0}}
\newcommand{\Lone}{\mathrm{L1}}
\newcommand{\Ltwo}{\mathrm{L2}}
\newcommand{\NT}{\textnormal{\textsc{not-testable}}}
\newcommand{\PROVED}{\textnormal{\textsc{proved}}}
\newcommand{\OPEN}{\textnormal{\textsc{open}}}
\newcommand{\e}{\mathrm e}
\newcolumntype{Y}{>{\raggedright\arraybackslash}X}

\title{\textbf{A Small-Polylogarithmic Determinant-Two}\\
\textbf{M\"obius Corridor: Uniform Actual-Core}\\
\textbf{Cancellation with Bounded Periodic Data}}
\author{Liang Wang\\
School of Artificial Intelligence and Automation,\\
Huazhong University of Science and Technology, Wuhan 430074, P.R. China\\
\texttt{wangliang.f@gmail.com}}
\date{July 2026}

\begin{document}
\maketitle

\begin{abstract}
We combine the general quantitative correlation theorem of Tao and
Ter\"av\"ainen with two exact interfaces.  TPC-148 constructs
multiplicative quotient lifts \(G_c\) satisfying
\(G_c(cm)=\mu(m)\), and TPC-147 reassembles a bounded periodic
multiplier without paying its number of residue classes.  The result
is a uniform theorem on the literal determinant-two M\"obius
periodic core.

There are absolute constants \(\eta_0,\kappa_0>0\) such that, for
each sufficiently large ambient \(X\), outside a subset of
\([\sqrt X,X]\) of normalized
logarithmic measure \(O((\log X)^{-\kappa_0})\), one has
\[
 \frac{as}{N}
 \left|
 \sum_{N<ad+asz\le2N}
 \mu(d+sz)\mu(u+az)\rho(z)
 \right|
 \ll\|\rho\|_\infty(\log X)^{-\kappa_0}
\]
uniformly for coprime odd \(a,s\), determinant \(su-ad=2\), and
bounded periodic \(\rho\) of period \(R\), provided
\(asR\le(\log X)^{\eta_0}\).  The exceptional union is taken only
over the distinct pairs \((a,s)\); the source theorem is already
uniform in residues and allowed moduli.

This is a proved \(\Lone\) actual-core theorem.  The current frontier
archive has not supplied its occurrence lift, and the theorem does
not include a nonperiodic physical multiplier, a generic phase,
arbitrary intervals, all prefixes or a four-point kernel.  Hence it
is not positive \(\Ltwo\), has zero fixed-\(X\)-power exponent and
does not approach a prime-pair conclusion by itself.
\end{abstract}

\section{Inputs and notation}

Let
\begin{equation}
 D(z)=d+sz,\qquad V(z)=u+az,\qquad su-ad=2,
\label{eq:fiber}
\end{equation}
where \(a,s\) are positive, coprime and odd.  These are the literal
determinant-two forms from the fixed-shift carrier
\citep{WangTPC127}.  Define
\begin{equation}
 Q_{\rm fib}=as,\qquad t(z)=aD(z)=ad+Q_{\rm fib}z.
\label{eq:t}
\end{equation}
Then \(t(z)+2=sV(z)\).

TPC-148 constructs one-bounded multiplicative functions \(G_c\)
satisfying
\begin{equation}
 G_c(cm)=\mu(m)
\label{eq:quotient}
\end{equation}
and proves, uniformly for small-polylogarithmic \(c\),
\begin{equation}
 \exp M(G_c;X^2,\log^{1/125}X)
 \gg(\log X)^\alpha
\label{eq:nonpret}
\end{equation}
for some absolute \(\alpha>0\)
\citep{WangTPC148,MatomakiRadziwillTao2015}.  Consequently
\begin{equation}
 \mu(D(z))\mu(V(z))=G_a(t(z))G_s(t(z)+2).
\label{eq:core}
\end{equation}

TPC-147 shows that a bounded period-\(R\) multiplier on the
\(Q_{\rm fib}\)-progression may be resolved modulo
\(Q_{\rm fib}R\) without an \(R\)-census loss, provided this total
modulus lies in the quantitative source envelope
\citep{WangTPC147}.

\section{An individual-pair corridor}

We first retain one fixed pair \((a,s)\) while allowing all
intercepts, periods and periodic values.

\begin{proposition}[One quotient-pair exceptional set]
\label{prop:pair}
There is an absolute \(\beta>0\) such that the following holds for
all sufficiently large \(X\).  For every coprime odd pair \((a,s)\)
with
\[
 as\le(\log X)^\beta,
\]
there is a set
\(\mathcal E_{a,s;X}\subset[\sqrt X,X]\) satisfying
\begin{equation}
 \frac1{\log X}\int_{\mathcal E_{a,s;X}}\frac{dt}{t}
 \ll(\log X)^{-\beta}.
\label{eq:pair-exception}
\end{equation}
If
\[
 N\in[\sqrt X,X]\setminus\mathcal E_{a,s;X},
 \qquad asR\le(\log X)^\beta,
\]
then, uniformly in all integral \(d,u\) obeying \eqref{eq:fiber}
and all period-\(R\) functions \(\rho\),
\begin{equation}
 \frac{as}{N}
 \left|
 \sum_{\substack{z\in\mathbb Z\\N<ad+asz\le2N}}
 \mu(D(z))\mu(V(z))\rho(z)
 \right|
 \ll\|\rho\|_\infty(\log X)^{-\beta}.
\label{eq:pair-bound}
\end{equation}
\end{proposition}

\begin{proof}
Use \eqref{eq:nonpret} with
\(\mathcal L=(\log X)^\alpha\), decreasing \(\alpha\) if needed.
The nonpretentious branch of
\citet[Theorem~3.1]{TaoTeravainen2026} then gives an absolute source
exponent \(c>0\), an exceptional set of normalized measure
\(\ll\mathcal L^{-c}\), and a correlation estimate with the same
saving for all moduli \(W\le\mathcal L^c\).  Put
\(\beta=\alpha c\), decreasing it harmlessly to absorb source
constants.

Invoke the source with
\[
 g_1=G_a,\quad g_2=G_s,\quad h_1=0,\quad h_2=2.
\]
On \(t\equiv ad\pmod{as}\), equation \eqref{eq:core} is exact.
TPC-147 reassembles \(\rho\) by using the refined modulus
\(W=asR\); the period count cancels against the density of each
refined progression.  The source exceptional set is uniform in
\(W\), its residue and the two shifts, so it is independent of
\(d,u,R,\rho\).  This proves \eqref{eq:pair-bound}.
\end{proof}

\begin{remark}[Why the intercepts have no source-height cost]
The source variable is \(t\), and \(ad\) enters only through its
representative modulo \(asR\).  Theorem~3.1 is uniform in that
residue.  This differs from calling the affine Liouville special
case, where all four affine coefficients and constants must lie in
the small-polylogarithmic envelope.
\end{remark}

\section{One set for all small pairs}

The set in \cref{prop:pair} may depend on \((a,s)\).  We now pay this
dependence exactly once.

\begin{lemma}[Unique quotient-pair census]
\label{lem:pairs}
For \(Q\ge2\),
\begin{equation}
 \#\{(a,s)\in\mathbb N^2:as\le Q\}
 =
 \sum_{a\le Q}\left\lfloor\frac Qa\right\rfloor
 \le Q(1+\log Q).
\label{eq:pairs}
\end{equation}
The same bound holds after imposing coprimality and oddness.
\end{lemma}

\begin{proof}
Drop the restrictions and use
\(\sum_{a\le Q}a^{-1}\le1+\log Q\).
\end{proof}

\begin{theorem}[Uniform determinant-two M\"obius-periodic corridor]
\label{thm:main}
There are absolute constants \(\eta_0,\kappa_0>0\) such that for
every sufficiently large \(X\) there is a set
\(\mathcal E_X^\star\subset[\sqrt X,X]\) with
\begin{equation}
 \boxed{
 \frac1{\log X}\int_{\mathcal E_X^\star}\frac{dt}{t}
 \ll(\log X)^{-\kappa_0}}
\label{eq:main-exception}
\end{equation}
for which the following is true.  Uniformly over all data satisfying
\begin{align}
 &a,s,R\in\mathbb N,\qquad d,u\in\mathbb Z,\qquad
 (a,s)=1,\qquad as\ {\rm odd},\qquad su-ad=2,\notag\\
 &asR\le(\log X)^{\eta_0},\qquad
 N\in[\sqrt X,X]\setminus\mathcal E_X^\star,
\label{eq:envelope}
\end{align}
and all \(\rho:\mathbb Z/R\mathbb Z\to\mathbb C\),
\begin{equation}
 \boxed{
 \frac{as}{N}
 \left|
 \sum_{\substack{z\in\mathbb Z\\N<ad+asz\le2N}}
 \mu(d+sz)\mu(u+az)\rho(z)
 \right|
 \ll
 \|\rho\|_\infty(\log X)^{-\kappa_0}.}
\label{eq:main}
\end{equation}
\end{theorem}

\begin{proof}
Let \(\beta\) be as in \cref{prop:pair}, and choose
\[
 0<\eta_0<\beta/4.
\]
Take the union of \(\mathcal E_{a,s;X}\) over all ordered pairs with
\(as\le(\log X)^{\eta_0}\).  By
\cref{lem:pairs,eq:pair-exception}, its normalized logarithmic
measure is
\[
 \ll
 (\log X)^{\eta_0+o(1)}(\log X)^{-\beta}
 =
 (\log X)^{-(\beta-\eta_0)+o(1)}.
\]
Choose, for example, \(\kappa_0=\beta/2\), after increasing the
lower threshold for \(X\).  Outside this union, the estimate
\eqref{eq:pair-bound} holds for every pair.  Its
\((\log X)^{-\beta}\) right side is stronger than
\((\log X)^{-\kappa_0}\).

No union over \(d,u,R\), residue classes or values of \(\rho\) is
made: for fixed \(G_a,G_s\), the source exceptional set is already
uniform in every allowed modulus and residue, and TPC-147 uses this
same set for the full periodic reassembly.
\end{proof}

\section{Squarefree and periodic reassembly audit}

\begin{proposition}[Where the previous CRT costs went]
\label{prop:costs}
Within \cref{thm:main}:
\begin{enumerate}[label=(\roman*),leftmargin=2.2em]
\item the two full M\"obius factors are represented exactly by
  \(G_a,G_s\), so there is no squarefree cutoff and no squarefree
  tail;
\item the \(R\) periodic residue classes have total native mass
  \(R\cdot N/(asR)=N/(as)\), so there is no \(R\)-census exponent;
\item the exceptional-set union is charged only to distinct
  multiplicative-function pairs \((G_a,G_s)\), equivalently to
  distinct \((a,s)\) in the declared range.
\end{enumerate}
\end{proposition}

\begin{proof}
Part (i) is \eqref{eq:core}, part (ii) is the TPC-147 periodic
reassembly identity, and part (iii) is the construction of
\(\mathcal E_X^\star\) in \cref{thm:main}.
\end{proof}

Products of literal bounded periodic factors may be absorbed in
\(\rho\) if their exact combined period keeps \(asR\) inside
\eqref{eq:envelope}.  This can include a proved periodic coprimality
mask, a fixed-period factor, or a rational phase whose exact
denominator is small enough.  It does not include a generic real
phase or a nonperiodic physical multiplier.

\section{Actual-core scope versus the actual archive}

The adjective ``actual-core'' means that the two M\"obius forms and
the determinant \(2\) are literal; it does not assert that every
frontier path has been lifted to such a record.  The current
TPC-143--146 interface leaves the occurrence-level fields
\[
 (a,s,d,u),\quad as,\quad\text{residue},\quad\text{ordered interval},
 \quad R,\quad\text{masks and coefficient mass}
\]
as \texttt{REQUIRED\_MISSING}.  Therefore the present frontier
consumption remains
\[
 \boxed{\mathsf{H1.frontier\_occurrence\_lift}=\NT.}
\]
The machine certificate accepts these missing obligations; it does
not manufacture numerical affine data from the finite cut sample.

Even after an occurrence lift is supplied, \cref{thm:main} has the
source interval
\[
 N<ad+asz\le2N
\]
at a nonexceptional scale.  A different interval must come with an
exact source-interval decomposition and an exceptional-set return.
All initial prefixes are not automatic.

\section{Claim boundary}

\begin{center}
\small
\begin{tabularx}{\textwidth}{@{}p{0.50\textwidth}Y@{}}
\toprule
Statement & Level and status\\
\midrule
Exact quotient-M\"obius determinant-two core &
\(\Lzero\), \(\PROVED\).\\
Uniform small-polylog periodic core theorem
\eqref{eq:main} &
Actual-core \(\Lone\), \(\PROVED\).\\
Full squarefree return in this two-point core &
Exact; no truncation tail.\\
Bounded periodic return inside \(asR\) envelope &
\(\Lone\), \(\PROVED\), no \(R\) loss.\\
Complete current frontier occurrence lift &
\(\NT\); required fields are absent.\\
Nonperiodic physical weight or generic phase &
\(\OPEN\).\\
All prefixes or four-point Fej\'er input &
\(\OPEN\).\\
Positive \(\Ltwo\), fixed \(X\)-power, \(1/400\) &
Not proved.\\
Prime-pair or twin-prime theorem & Not proved.\\
\bottomrule
\end{tabularx}
\end{center}

The theorem is forward arithmetic progress: it replaces a
squarefree-CRT shadow by a complete M\"obius core on the correct
fixed-two geometry.  Its next legal consumer must solve the local
exceptional-window and deterministic-prefix return without
promoting logarithmic cancellation to an \(X\)-power.

\begingroup
\footnotesize
\raggedright
\bibliographystyle{plainnat}
\bibliography{references}
\endgroup

\end{document}
